\section{Area, spherical excess, solid angle, and polygons}
\label{sec:area}

\subsection{Constructive spherical area}

A statement about area requires more than the algebraic triangle certificate,
but it still admits a primitive-recursive finite-precision interpretation.
Use the equal-area coordinates
\begin{equation}
 \Phi(\lambda,z)=
 \bigl(\sqrt{1-z^2}\cos\lambda,
       \sqrt{1-z^2}\sin\lambda,z\bigr),
 \quad -\pi\le\lambda\le\pi,\quad -1\le z\le1.            \tag{10.1}
\end{equation}
Away from the two pole lines,
\begin{equation}
 \norm{\partial_\lambda\Phi\crossp\partial_z\Phi}=1.       \tag{10.2}
\end{equation}
Thus surface area is ordinary rectangle area \(d\lambda\,dz\) in this chart.
All coordinate functions have the constructive names from Section~\ref{sec:ipra}.

For a polygon bounded by \(m\) great-circle arcs, subdivide the
\((\lambda,z)\)-rectangle into a dyadic grid.  A cell whose interval evaluation
has a fixed sign against every boundary normal is decidably inside or outside.
The undecided cells lie in a union of strips around the boundary great circles.
The strip of angular half-width \(\delta\) around one great circle has area
\begin{equation}
 4\pi\sin\delta\le4\pi\delta.                              \tag{10.3}
\end{equation}
Therefore the gap between lower and upper cell sums is at most
\(4\pi m\delta\), plus the explicitly bounded coordinate-roundoff margin.
Choosing a dyadic \(\delta\) from the requested output precision gives a
primitive-recursive modulus.  This constructs a rapid area name.  Cellwise
counting proves finite additivity; boundary strips vanish with the same
modulus.  No appeal to a nonconstructive measure-extension theorem is needed.

There is an even simpler computation for a lune.  Choose the polar axis of its
two bounding great circles as the \(z\)-axis.  A lune of angle \(\alpha\)
corresponds to the rectangle
\(0\le\lambda\le\alpha\), \(-1\le z\le1\).  Equations
(10.1)--(10.2) give
\begin{equation}
 \Area(\text{lune of angle }\alpha)=2\alpha.               \tag{10.4}
\end{equation}
At finite precision the Riemann sum is exact because the area density is the
constant \(1\).

\subsection{Girard's theorem}

\begin{theorem}[Girard's spherical excess formula]
\label{thm:girard}
For a proper spherical triangle on the unit sphere,
\begin{equation}
 \Area(ABC)=E:=A+B+C-\pi.                                  \tag{10.5}
\end{equation}
On a sphere of radius \(R\), the area is \(R^2E\).
\end{theorem}

\begin{proof}
The three side great circles partition the sphere into eight spherical
triangles, paired antipodally.  Let the area of the original triangle be
\(T\), and choose one triangle from each of the other three antipodal pairs,
with areas \(X,Y,Z\), so that these four regions form a hemisphere.  Finite
additivity gives
\begin{equation}
 T+X+Y+Z=2\pi.                                              \tag{10.6}
\end{equation}
The interior-angle lune at \(A\) consists of the original triangle together
with one of the regions \(X,Y,Z\); after relabeling take its area to be
\(T+X\).  Similarly the \(B\)- and \(C\)-lunes have areas \(T+Y\) and
\(T+Z\).  By (10.4), their total area is
\begin{equation}
 2A+2B+2C=3T+X+Y+Z.                                        \tag{10.7}
\end{equation}
Use (10.6) to replace \(X+Y+Z\) by \(2\pi-T\).  Then
\[
 2(A+B+C)=2T+2\pi,
\]
so \(T=A+B+C-\pi\).  The proof is a finite dissection identity plus the exact
lune calculation.  Both are valid for the rapid area names just constructed.
Scaling a unit-sphere patch by \(R\) multiplies every area by \(R^2\).
\end{proof}

\subsection{An algebraic solid-angle formula}

The area can also be recovered directly from the vector certificate.  Put
\begin{equation}
 D=1+x+y+z,
 \qquad
 Q=4\cos\frac a2\cos\frac b2\cos\frac c2
   =\sqrt{2(1+x)(1+y)(1+z)}.                               \tag{10.8}
\end{equation}
The polynomial identity
\begin{equation}
 D^2+\tau^2=2(1+x)(1+y)(1+z)=Q^2                          \tag{10.9}
\end{equation}
uses only the Gram relation.

\begin{theorem}[Oriented solid-angle formula]
\label{thm:solid-angle}
For the positive orientation,
\begin{align}
 \cos\frac E2&=\frac{D}{Q},                                \tag{10.10}\\
 \sin\frac E2&=\frac{\tau}{Q},                             \tag{10.11}\\
 E&=2\atanTwo(\tau,D).                                     \tag{10.12}
\end{align}
If \(D\apart0\), then
\begin{equation}
 \tan\frac E2=\frac{\tau}{1+x+y+z}.                       \tag{10.13}
\end{equation}
\end{theorem}

\begin{proof}
Since \(E/2=P-\pi/2\), the addition formulas give
\(\cos(E/2)=\sin P\) and \(\sin(E/2)=-\cos P\).  Expand
\(\sin P\) and \(\cos P\) by successive half-angle addition.  Substitute the
half-angle formulas of Section~\ref{sec:half-formulas}, clear the positive
factor \(4\cos(a/2)\cos(b/2)\cos(c/2)\), and use (4.16).  The two numerators
reduce respectively to \(D\) and \(\tau\).  This is a finite polynomial
normalization after the certified square-root equations are inserted.
Equation (10.9) verifies that the displayed pair is a unit pair.  Since
\(0<E<2\pi\) and \(\tau>0\), the constructive two-chart \(\atanTwo\) operation
selects exactly \(E/2\), proving (10.12).  Guarded division yields (10.13).
\end{proof}

Formula (10.12) is often the most stable computational form of Girard's
excess.  It avoids subtracting nearly equal angles for a small triangle.

\subsection{L'Huilier's formula}

\begin{corollary}[L'Huilier]
\label{cor:lhuillier}
Let \(p=(a+b+c)/2\).  Then
\begin{equation}
 \tan^2\frac E4=
 \tan\frac p2
 \tan\frac{p-a}{2}
 \tan\frac{p-b}{2}
 \tan\frac{p-c}{2},                                       \tag{10.14}
\end{equation}
and the positive square root gives \(\tan(E/4)\).
\end{corollary}

\begin{proof}
From (10.10),
\begin{equation}
 \tan^2\frac E4=\frac{1-\cos(E/2)}{1+\cos(E/2)}
 =\frac{Q-D}{Q+D}.                                         \tag{10.15}
\end{equation}
Write each tangent on the right of (10.14) as a sine divided by a cosine.
Apply product-to-sum twice to the four numerator factors and twice to the four
denominator factors.  After inserting
\(Q=4\cos(a/2)\cos(b/2)\cos(c/2)\), the resulting numerator is
\(Q-D\) and the denominator is \(Q+D\).  Equivalently, cross multiplication
reduces to (10.9) and the three identities
\(2\sin r\sin s=\cos(r-s)-\cos(r+s)\).  All factors in (10.14) are positive
for a proper triangle, so nonnegative-root uniqueness fixes the sign.
\end{proof}

\subsection{Spherical polygons}

\begin{theorem}[Polygonal excess]
\label{thm:polygon-excess}
Let a positively oriented simple geodesic \(n\)-gon on the unit sphere have
interior angles \(\alpha_1,\ldots,\alpha_n\), and suppose a certified
noncrossing triangulation lies in the polygon.  Then
\begin{equation}
 \Area(P)=\sum_{j=1}^n\alpha_j-(n-2)\pi.                   \tag{10.16}
\end{equation}
On a radius-\(R\) sphere, multiply by \(R^2\).
\end{theorem}

\begin{proof}
A triangulation has \(n-2\) triangles.  Finite additivity says the polygonal
area is the sum of their areas.  The triangle angles partition the polygon's
interior angles, including at vertices incident to several diagonals.  Applying
Girard to each triangle therefore gives
\[
 \Area(P)
 =\sum_{j=1}^n\alpha_j-(n-2)\pi.
\]
The entire argument is finite.  The triangulation and its noncrossing/orientation
signs are part of the arithmetic certificate, so no choice of a diagonal is
made inside \(\ipra\).
\end{proof}
